\section{Reduced densities of torus cuts}%
\label{sec:peps_torus_cut_densities}

This section first records the entropies of a density operator with flat
spectrum. It then computes the reduced density of the physical cut from the
common boundary factor, for a fixed boundary state and for the closures of the
original tensor.

\begin{lemma}[Entropy of a flat density operator]%
    \label{thm:peps_flat_density_entropy}
    \lean{Matrix.IsHermitian.eigenvalues_eq_zero_or_inv_of_mul_self_eq_inv_smul,
        vonNeumannEntropy_of_mul_self_eq_inv_smul}
    \leanok
    Let $\rho$ be a finite-dimensional positive semidefinite matrix of
    trace one. As in the spectral step of the proof of
    \cite[Theorem~6.9]{Schuch2010PEPS}, if $\rho^2=r^{-1}\rho$ for a positive real number $r$, then
    $S(\rho)=\log r$.
\end{lemma}
\begin{proof}\leanok
    For each eigenvalue $\lambda$, the identity gives
    $\lambda(\lambda-r^{-1})=0$. In either case
    $-\lambda\log\lambda=\lambda\log r$, with the value at zero
    interpreted continuously. Summing and using
    $\sum\lambda=\operatorname{tr}\rho=1$ proves $S(\rho)=\log r$.
\end{proof}

\begin{definition}[Finite-order R\'enyi entropy]%
    \label{def:peps_renyi_entropy}
    \lean{renyiEntropy}\leanok
    For a positive density $\rho$ and $\alpha>0$, $\alpha\ne1$, define
    \begin{align}
        S_\alpha(\rho)=(1-\alpha)^{-1}
        \log\operatorname{tr}(\rho^\alpha),
    \end{align}
    where the real power is defined spectrally. Set $S_1(\rho)=S(\rho)$
    and $S_0(\rho)=\log\operatorname{rank}\rho$.
\end{definition}

\begin{lemma}[All finite R\'enyi orders of a flat density]%
    \label{thm:peps_flat_density_renyi_entropy}
    \lean{trace_cfc_rpow_of_mul_self_eq_inv_smul,
        renyiEntropy_of_mul_self_eq_inv_smul,
        renyiEntropy_of_mul_self_eq_inv_smul_of_rank}
    \leanok
    \uses{def:peps_renyi_entropy}
    Let $\rho\geq0$ have trace one and satisfy $\rho^2=r^{-1}\rho$, with
    $r>0$. For every $\alpha>0$,
    $\operatorname{tr}(\rho^\alpha)=r^{1-\alpha}$ and
    $S_\alpha(\rho)=\log r$. If $\operatorname{rank}\rho=r$, the same
    entropy formula also holds at order zero.
\end{lemma}
\begin{proof}\leanok
    \uses{thm:peps_flat_density_entropy}
    Every eigenvalue is zero or $r^{-1}$. For positive order,
    $\lambda^\alpha=\lambda r^{1-\alpha}$. Summing uses trace one;
    taking logarithms proves the entropy formula away from order one.
    The von Neumann and rank conventions give the remaining orders.
\end{proof}

\begin{definition}[Fixed boundary embedding and density]%
    \label{def:peps_fixed_boundary_factor_density}
    \lean{TNLean.PEPS.boundaryAncillaMap,
        TNLean.PEPS.fixedBoundaryFactorDensity,
        TNLean.PEPS.boundaryFactorSchmidtMatrix}
    \leanok
    Let $D$ be a nonempty finite set and $\phi\in\mathbb C^A$.
    Define $J_\phi:\mathbb C^D\to\mathbb C^{D\times A}$ by
    $(J_\phi)_{(d,a),s}=\mathbf 1_{\{d=s\}}\phi(a)$, and set
    \begin{align}
        \rho_\phi=|D|^{-1}J_\phi J_\phi^\dagger,
        \qquad K_{\phi,\chi}=|D|^{-1/2}J_\phi J_\chi^{\mathsf T}.
    \end{align}
\end{definition}

\begin{lemma}[Reduced matrix of a fixed boundary factor]%
    \label{thm:peps_fixed_boundary_factor_density}
    \lean{TNLean.PEPS.boundaryAncillaMap_conjTranspose_mul,
        TNLean.PEPS.fixedBoundaryFactorDensity_properties,
        TNLean.PEPS.vonNeumannEntropy_fixedBoundaryFactorDensity,
        TNLean.PEPS.boundaryFactorSchmidtMatrix_apply,
        TNLean.PEPS.boundaryFactorSchmidtMatrix_mul_conjTranspose,
        TNLean.PEPS.mul_unitary_transpose_mul_conjTranspose,
        TNLean.PEPS.boundaryFactorSchmidtMatrix_submatrix_mul_conjTranspose,
        TNLean.PEPS.fixedBoundaryFactorDensity_submatrix_properties,
        TNLean.PEPS.normalized_reducedMatrix_eq_of_scalar}
    \leanok
    \uses{def:peps_fixed_boundary_factor_density}
    If $\|\phi\|=1$, then $J_\phi$ is an isometry and $\rho_\phi$ is
    a positive density of trace one and rank $|D|$, satisfying
    $\rho_\phi^2=|D|^{-1}\rho_\phi$. In addition,
    \begin{align}
        K_{\phi,\chi}K_{\phi,\chi}^\dagger
        =\|\chi\|^2\rho_\phi.
    \end{align}
    These identities persist under bijective changes of physical coordinates.
    A unitary acting on the complementary indices leaves the reduced matrix
    unchanged. If a nonzero coefficient matrix has reduced matrix
    $a\rho_\phi$, its trace is positive and trace normalization gives
    exactly $\rho_\phi$.
\end{lemma}
\begin{proof}\leanok
    \uses{thm:peps_flat_density_entropy}
    The Gram matrix is $J_\phi^\dagger J_\phi=\|\phi\|^2I$.
    Matrix multiplication gives the reduced identity. Cyclicity of trace,
    the isometry rank identity, and one further multiplication give the
    density properties. Complementary unitarity cancels between a matrix
    and its adjoint; positivity and nonvanishing give positive trace.
\end{proof}

\begin{definition}[Fixed regular reference-boundary density]%
    \label{def:peps_regular_reference_boundary_density}
    \lean{TNLean.PEPS.RegularReferenceRegionLabels,
        TNLean.PEPS.regularReferenceRegionLabelsFintype,
        TNLean.PEPS.regularRegionReferenceCoordinatesEquiv,
        TNLean.PEPS.regularReferenceBoundaryDensity}
    \leanok
    \uses{def:peps_regular_region_coordinates,
        thm:peps_controlled_projector_cut_boundary_factor,
        def:peps_fixed_boundary_factor_density}
    Number the $n+1$ crossing bonds. The spanning-tree physical
    coordinates identify the region rows with $(d,a)$, where $d\in G^n$
    is the relative boundary label and $a$ consists of its reference label,
    internal references, normalized vertex labels, and cycle residuals.
    Pull back $\rho_{\phi_R}$ to the native half-edge physical rows, with
    the fixed normalized ancillary vector $\phi_R$ of
    Lemma~\ref{thm:peps_controlled_projector_cut_boundary_factor}.
    Denote this density by $\rho_R^0$.
\end{definition}

\begin{lemma}[Spectrum of the fixed regular reference density]%
    \label{thm:peps_regular_reference_boundary_density}
    \lean{TNLean.PEPS.star_dotProduct_regularReferenceRegionAncilla,
        TNLean.PEPS.regularReferenceBoundaryDensity_properties,
        TNLean.PEPS.vonNeumannEntropy_regularReferenceBoundaryDensity,
        TNLean.PEPS.renyiEntropy_regularReferenceBoundaryDensity}
    \leanok
    \uses{def:peps_regular_reference_boundary_density}
    The density $\rho_R^0$ is positive, has trace one and rank $|G|^n$, and
    satisfies $(\rho_R^0)^2=|G|^{-n}\rho_R^0$. Its von Neumann entropy and
    every nonnegative finite-order R\'enyi entropy are $n\log|G|$.
\end{lemma}
\begin{proof}\leanok
    \uses{thm:peps_fixed_boundary_factor_density, thm:peps_flat_density_renyi_entropy}
    Since the fixed ancillary vector has norm one, applying the fixed-factor density
    identities and pulling them back through the coordinate bijection yields the result.
\end{proof}

\begin{lemma}[Common reduced density of canonical closure sums]%
    \label{thm:peps_torus_controlled_boundary_density}
    \lean{TNLean.PEPS.torusControlledCut_reducedMatrix_eq_trace_smul,
        TNLean.PEPS.torusControlledCut_normalized_reducedMatrix,
        TNLean.PEPS.torusControlledCut_normalized_renyiEntropy}
    \leanok
    \uses{thm:peps_torus_controlled_boundary_factor}
    Retain the simple connectedness, fixed spanning trees, crossing
    numbering, and fixed complementary winding loops of
    Lemma~\ref{thm:peps_torus_controlled_boundary_factor}. For a finite
    commuting closure family, let $C_\lambda$ be its coherent
    canonical cut matrix, before the complementary controlled unitary.
    Then the fixed density $\rho_R^0$, chosen before the sectors and
    coefficients, satisfies
    \begin{align}
        C_\lambda C_\lambda^\dagger
        =\operatorname{tr}(C_\lambda C_\lambda^\dagger)\rho_R^0.
    \end{align}
    If $C_\lambda\ne0$, its trace is positive and its normalized reduced
    density equals $\rho_R^0$. Consequently its rank is $|G|^n$, its
    nonzero spectrum is flat, and its von Neumann and all nonnegative
    finite-order R\'enyi entropies are $n\log|G|$.

    This is a canonical-cut consequence with given fixed complementary
    winding loops. Their existence for the paper's arbitrary disk geometry
    and transfer to the original site tensors remain separate steps;
    see \cite{gap:rmp_peps_quantum_double_g_isometry}. It is not an
    unrestricted formalization of \cite[Theorem~6.9]{Schuch2010PEPS}.
\end{lemma}
\begin{proof}\leanok
    \uses{thm:peps_regular_reference_boundary_density}
    The controlled cut is $K_{\phi_R,\chi_\lambda}$ in fixed physical
    coordinates. Its reduced matrix is $\|\chi_\lambda\|^2\rho_R^0$.
    Complementary unitarity gives the same reduced matrix before control.
    Since $\operatorname{tr}\rho_R^0=1$, the scalar is precisely its trace.
    Nonvanishing then permits normalization; the fixed spectrum gives all
    stated entropies.
\end{proof}

\begin{lemma}[One complementary unitary from fixed winding data]%
    \label{thm:peps_torus_complement_disentangling}
    \lean{TNLean.PEPS.exists_unitary_torusComplementDisentangling_of_winding_eq}
    \leanok
    \uses{thm:peps_torus_complement_cycle_words, thm:peps_cycle_controlled_boundary}
    Choose the two rooted trees, reference crossing bond, and fixed
    complementary root loops with matching winding from
    Lemma~\ref{thm:peps_torus_complement_cycle_words}. There exists one
    unitary $U$ on the original complementary half-edge physical space,
    chosen independently of the commuting closure labels $g,h$.
    Write $S=R^c$, let $B_S^{g,h}$ be the canonical complementary
    block, and let $F_R^{g,h},F_S^{g,h}$ be the boundary maps.
    Reread each crossing label on $S$ under the same native edge, and put
    $m_f^{g,h}=(F_S^{g,h}\circ(F_R^{g,h})^{-1})(\mathbf 1)_f$.
    For every commuting pair and every complementary physical coordinate
    $(y,a,r,z)$, the coefficient formula is
    \begin{align}
         & (U B_S^{g,h})\bigl(\alpha(y,a,r,z),
        (F_R^{g,h})^{-1}\theta\bigr)           \\
         & \qquad=|G|^{-|S|}\sum_{x\in G}
        \mathbf 1_{\{y_f=xm_{f_0}^{g,h}\theta_f\ \text{for every }f\}}
        \mathbf 1_{\{z=x\omega_S^{g,h}x^{-1}\}}.
        \label{eq:peps_complement_disentangled_coefficient}
    \end{align}
    Thus the relative crossing multipliers are removed by one common
    complementary operation, while the reference multiplier and
    cycle constraint remain explicit. This is an auxiliary form of the
    disentangling calculation in
    \cite[Theorem~6.9]{Schuch2010PEPS}, local source lines 1935--1990.
    Lemma~\ref{thm:peps_torus_simply_connected_boundary_loops}
    gives these winding loops from simple connectedness.
\end{lemma}

\begin{proof}\leanok
    \uses{thm:peps_cycle_controlled_relative_walks}
    At cycle coordinates $z$, translating the boundary coordinate $y_f$
    by the inverse of the evaluation of the fixed complementary word
    associated with $f$ while leaving the other coordinates unchanged
    defines a permutation of the original half-edge physical basis, and hence
    a unitary independent of $g,h$. Matching winding derives the
    relative word identity for each commuting closure pair. Applying
    the exact controlled coefficient formula then cancels the relative
    multipliers, retaining the reference and cycle factors of
    (\ref{eq:peps_complement_disentangled_coefficient}).
\end{proof}

\begin{lemma}[Local invariant support of walk-controlled operations]%
    \label{thm:peps_cycle_controlled_local_support}
    \lean{TNLean.PEPS.regularRegionCoordinatesEquiv_symm_coordinateTranslation,
        TNLean.PEPS.regularRegionCoordinateTranslation,
        TNLean.PEPS.regularCycleControlledBoundary_coordinateTranslation,
        TNLean.PEPS.regularRegionVertexTranslationMatrix,
        TNLean.PEPS.regularCycleControlledBoundaryMatrix_commute_vertexTranslation,
        TNLean.PEPS.regularRegionVertexTranslationMatrix_apply,
        TNLean.PEPS.regionPhysicalProductMatrix_regularLegProjector_eq_sum_vertexTranslation,
        TNLean.PEPS.regularCycleControlledBoundaryMatrix_commute_localProjector,
        TNLean.PEPS.regularCycleControlledBoundaryMatrix_maps_localProjector_range,
        TNLean.PEPS.exists_unitary_regionPhysicalProductMatrix_mul_cycleWords}
    \leanok
    \uses{def:peps_regular_region_coordinates, thm:peps_regular_region_cycle_holonomy, thm:peps_cycle_controlled_boundary, thm:peps_region_physical_maps}
    Fix a rooted spanning tree in a finite region. Independent vertex
    translations $\lambda_v$ of its physical half-edge labels
    transform the spanning-tree coordinates by
    \begin{align}
        (y,a,r,z) & \longmapsto
        \bigl(\lambda_oy,\lambda_oa,
        (\lambda_vr_v\lambda_o^{-1})_v,
        \lambda_oz\lambda_o^{-1}\bigr).
        \label{eq:peps_local_translation_coordinates}
    \end{align}
    Reconstructing these coordinates gives precisely the original
    half-edge labels multiplied on the left by $\lambda_v$ at vertex $v$.
    For a fixed family of region walks, let $U$ be the native
    boundary permutation controlled by their constructed cycle words.
    This matrix commutes with every independent vertex translation.
    Their finite average is the product local invariant projector
    \begin{align}
        P & =\bigotimes_{v\in R}P_v
        =|G|^{-|R|}\sum_{\lambda\in G^R}L_\lambda,
          & UP                      & =PU.
        \label{eq:peps_controlled_local_projector}
    \end{align}
    In particular $U$ maps the local invariant support
    $\operatorname{range}P$ into itself.

    If the original local physical maps $A_v$ are regular $G$-isometric,
    there exists a unitary $W$ on their physical region such that
    \begin{align}
        W\Bigl(\bigotimes_{v\in R}A_v\Bigr)
         & =\Bigl(\bigotimes_{v\in R}A_v\Bigr)U.
    \end{align}
    The fixed words and their physical implementation are chosen without
    closure labels. No equivariance premise on an arbitrary control
    function or surjectivity of the canonical projector is required.
    These are auxiliary physical support identities for the accessible
    complement calculation in \cite[\S6.2--6.3]{Schuch2010PEPS},
    local source lines 1765--1920 and 1935--1990.
\end{lemma}

\begin{proof}\leanok
    \uses{thm:peps_regular_physical_unitary_transport}
    Applying the native reconstruction formulas at the two internal endpoints
    and the region endpoint of each crossing edge causes the vertex coordinate
    in (\ref{eq:peps_local_translation_coordinates}) to cancel its root inverse against the transformed reference
    or boundary label, giving the independent physical translation.
    Simultaneous conjugation conjugates evaluation of every walk
    word, so the controlled boundary translation commutes with this action.
    Taking the product of the local finite averages identifies their
    independent choices with $\lambda\in G^R$, yielding the product
    projector and commutation (\ref{eq:peps_controlled_local_projector}). This also proves range preservation.
    The supported physical unitary implementation theorem then gives $W$ on the
    original region, including any unused physical directions.
\end{proof}

\begin{lemma}[Physical transfer of the canonical cut]%
    \label{thm:peps_regular_physical_cut_transfer}
    \lean{TNLean.PEPS.regionPhysicalProductMatrix_mul_regularProjectorTwistedRegionMatrix,
        TNLean.PEPS.regularProjectorTwistedCutMatrix,
        TNLean.PEPS.regionPhysicalProductMatrix_mul_regularProjectorTwistedCutMatrix,
        TNLean.PEPS.regularProjectorTwistedCutMatrix_mul_regionPhysicalProductMatrix_transpose,
        TNLean.PEPS.regularPhysicalCutMatrix_eq_physicalImage_regularProjectorTwistedCutMatrix,
        TNLean.PEPS.exists_positive_regularPhysicalCut_density_transfer}
    \leanok
    \uses{thm:peps_region_physical_maps}
    For arbitrary bond operators $u$, contract the two canonical
    projector regions over their common crossing labels:
    \begin{align}
        C_u(\alpha,\beta)
         & =\sum_\theta B_R^u(\alpha,\theta)
        B_{R^c}^u(\beta,\theta).
    \end{align}
    Here the complement reads the same crossing edge under its native
    boundary identification. Write $P_R=\bigotimes_{v\in R}P_v$ and
    $T_R=\bigotimes_{v\in R}A_v$ for the original local projector and
    physical maps. The physical coefficient matrix $M_u$ satisfies
    \begin{align}
        P_R C_u & =C_u,                        & C_uP_{R^c}^{\mathsf T} & =C_u, &
        M_u     & =T_R C_uT_{R^c}^{\mathsf T}.
    \end{align}
    If the local maps are regular $G$-isometric, choose their derived
    positive product Gram factors $c_R,c_{R^c}$. For every finite coherent
    sum $C=\sum_i\mu_iC_{u_i}$ and $M=\sum_i\mu_iM_{u_i}$, put
    $\tau_C=CC^\dagger$ and $\tau_M=MM^\dagger$. Then
    \begin{align}
        \tau_M & =c_{R^c}T_R\tau_CT_R^\dagger,       &
        \operatorname{tr}\tau_M
               & =c_Rc_{R^c}\operatorname{tr}\tau_C,   \\
        (\operatorname{tr}\tau_M)^{-1}\tau_M
               & =c_R^{-1}T_R
        \bigl((\operatorname{tr}\tau_C)^{-1}\tau_C\bigr)T_R^\dagger.
    \end{align}
    The inverse is read as zero when its argument is zero. For a nonzero
    state, the last identity is the transfer of normalized reduced density.
    This is an auxiliary physical accessibility calculation for
    \cite[Theorem~6.9]{Schuch2010PEPS}, local source lines 1765--1820 and
    1935--1990; no exterior routing or canonical factorization is assumed.
\end{lemma}

\begin{proof}\leanok
    \uses{thm:peps_regular_physical_unitary_transport, thm:peps_regular_twisted_region_accessibility}
    Expanding the half-edge contraction and summing the physical row
    against the product local projector replaces each site factor
    $P_v$ by $P_v^2=P_v$. The two reverse accessibility identities give
    $M_u=T_RC_uT_{R^c}^{\mathsf T}$ directly. These equalities persist
    under finite coherent sums. Inserting
    $T_{R^c}^\dagger T_{R^c}=c_{R^c}P_{R^c}$ gives the reduced matrix;
    cyclicity of trace and $P_RC=C$ give its trace. Cancelling the
    positive product factors gives the normalized identity.
\end{proof}

\begin{lemma}[Spectrum under a supported physical map]%
    \label{thm:peps_supported_physical_density}
    \lean{TNLean.PEPS.normalizedPhysicalImageDensity_properties,
        TNLean.PEPS.normalizedPhysicalImageDensity,
        TNLean.PEPS.vonNeumannEntropy_normalizedPhysicalImageDensity}
    \leanok
    \uses{thm:peps_regular_physical_cut_transfer}
    Let $\rho\geq0$ have trace one, and suppose a rectangular physical map
    $T$ and a positive scalar $c$ satisfy $T^\dagger T\rho=c\rho$.
    Define $\widehat\rho=c^{-1}T\rho T^\dagger$. Then
    $\widehat\rho\geq0$, $\operatorname{tr}\widehat\rho=1$, and
    $\operatorname{rank}\widehat\rho=\operatorname{rank}\rho$.
    If $\rho^2=r^{-1}\rho$, then
    \begin{align}
        \widehat\rho^2  & =r^{-1}\widehat\rho, &
        S(\widehat\rho) & =\log r.
    \end{align}
    This auxiliary finite-dimensional calculation expresses physical
    accessibility on the invariant support in
    \cite[Theorem~6.9]{Schuch2010PEPS}, local source lines 1765--1820 and
    2043--2072. It makes no assertion that $T$ is surjective.
\end{lemma}

\begin{proof}\leanok
    \uses{thm:peps_flat_density_entropy}
    Hermitian conjugation gives $\rho T^\dagger T=c\rho$.
    Positivity follows by conjugation and positive scaling; cyclicity of
    trace gives trace one. In addition,
    \begin{align}
        T^\dagger\widehat\rho T & =c\rho.
    \end{align}
    Applying the rank inequality to each physical image yields equality
    of ranks. Inserting $T^\dagger T$ between the two densities proves
    the flat-spectrum identity, which gives the stated entropy.
\end{proof}

\begin{lemma}[Common reduced density for original torus closures]%
    \label{thm:peps_torus_physical_cut_density}
    \lean{TNLean.PEPS.IsGIsometric.exists_torusPhysicalCut_density_transfer,
        TNLean.PEPS.IsGIsometric.exists_torusPhysicalCut_common_density_of_isSimplyConnected_and_winding}
    \leanok
    \uses{thm:peps_regular_physical_cut_transfer, thm:peps_supported_physical_density}
    Work on a native torus with both dimensions at least three, with an
    original regular $G$-isometric site map. Assume that the
    closed-cell realization of $R$ is simply connected. Choose rooted
    spanning trees on $R$ and its complement, a nonempty enumeration
    of the $b$ crossing bonds, and fixed complement-root loops whose
    winding equals the corresponding two-tree boundary routes.
    There exists one density $\rho_R$ on the original physical region,
    chosen before the closure pairs and their coherent coefficients.
    For every finite family of commuting pairs $(g_i,h_i)$ and every
    nonzero coherent closure superposition with cut matrix $M$,
    \begin{align}
        \frac{MM^\dagger}{\operatorname{tr}(MM^\dagger)} & =\rho_R,             &
        \rho_R                                           & \geq0,               & \operatorname{tr}\rho_R & =1, \\
        \operatorname{rank}\rho_R                        & =|G|^{b-1},          &
        \rho_R^2                                         & =|G|^{-(b-1)}\rho_R, &
        S(\rho_R)                                        & =(b-1)\log|G|.
    \end{align}
    For every nonnegative finite Rényi order $\alpha$, including order
    zero and the von Neumann value at order one,
    $S_\alpha(\rho_R)=(b-1)\log|G|$.
    More generally, the density of every coherent twisted cut
    is the supported physical image of its canonical density, with
    one positive normalization independent of the closure labels.
    The first conclusion requires no surjectivity of the original
    physical map, assumed block Gram, or tensor factorization.
    This is an auxiliary form of \cite[Theorem~6.9]{Schuch2010PEPS},
    local source lines 1935--1990 and 2027--2072. The fixed exterior
    winding loops are explicit hypotheses here;
    Lemma~\ref{thm:peps_torus_simply_connected_physical_density}
    derives them from the contiguous simply connected region.
\end{lemma}

\begin{proof}\leanok
    \uses{thm:peps_torus_controlled_boundary_density, thm:peps_flat_density_renyi_entropy}
    Deriving the incident-coordinate local isometries from the original
    four-leg isometry gives a physical product map with positive scalar
    Gram $T_R^\dagger T_R=c_RP_R$. The cut image shows that a
    nonzero original coherent cut has a nonzero canonical coherent cut.
    Its fixed normalized density $\rho_R^0$ is supported on $P_R$, so
    the physical density is the fixed matrix
    \begin{align}
        \rho_R & =c_R^{-1}T_R\rho_R^0T_R^\dagger.
    \end{align}
    The supported physical density theorem transfers positivity, trace,
    rank, and the flat spectrum. The spectral entropy calculation gives
    the stated boundary term at all nonnegative finite Rényi orders.
    Every physical map and normalization was
    chosen from the original site tensor before the coherent coefficients.
\end{proof}
